

\label{ch:introduction}\section{Introduction and Opening Remarks} 

\textcolor{red}{To-Do: Rewrite Intro, Background to push for design-for-lifetime security. Pull Related Works from Dissertation. Finalize Completed Work writeup (Sec.4) in Dissertation and copy/paste here. Add in that we generate the access control policies. Discuss redundancy and inter-component communication risks. Discuss original metric and why it was not sufficient in related works. Print results from existing Python code from prior submission. Max 12 pages}

Today, embedded systems and Cyber-Physical Systems (CPS) increasingly rely on Systems-on-Chip (SoCs). These specialized, often reconfigurable devices integrate diverse hardware components and software onto a single chip, offering enhanced performance, flexibility, and time-to-market. However, the growing complexity of these devices makes it impractical for companies to design the entire system in-house. Consequently, the integration of third-party components, many of which may not be fully trusted, poses significant security risks when operating unguarded in critical applications.

Existing security mitigation approaches, such as pre-deployment Trojan detection for Intellectual Property (IP) blocks, require extensive testing and may not eliminate all malicious elements. Traditional embedded systems and operational technology (OT) systems have largely relied on post-design activities like antivirus software and firewalls, which primarily aim to prevent malicious actors from accessing and disrupting system operation. In both cases, security is not considered an integral part of the design process. 

The target of this work is to design systems that allow trusted and untrusted components to operate alongside each other, equipped with mechanisms to protect system assets and prevent potential attacks and service disruptions. To achieve this, the base SoC architecture is expanded with hardware and software security primitives that strictly adhere to the zero-trust principle for asset protection. We present a novel design automation framework for secure SoCs, where security is an integral part of the system generation and optimization process. Given an SoC with specified security requirements, we propose an automated system generation approach that simultaneously optimizes system resources, power, and security objectives. Our method leverages Answer Set Programming (ASP) to express the complex multi-objective design problem in a declarative manner. This approach enables comprehensive design space exploration to obtain diverse sets of solutions with desirable properties. While design space exploration has been extensively researched, and qualitative solutions exist, the nature of security properties, often expressed as advanced logic, makes ASP an ideal candidate for the design space exploration of secure SoCs, surpassing the capabilities of traditional evolutionary algorithms, integer linear programming, or particle swarm optimization for this specific challenge.

We provide a novel formulation of secure SoC synthesis as a multi-objective optimization problem within the ASP framework, which is then solved using the Potsdam Answer Set Solving Collection (Potassco) \cite{gebser_potassco_2011}, described further in Section \ref{sec:problem}. The resulting solutions are formatted into interpretable graphs that system designers can utilize to determine their system's optimal configuration. Our design was validated through representative case studies and a synthetic benchmark, demonstrating the viability of our approach. The full ASP solver and Python code is provided in our attached GitHub. \footnote{https://github.com/smartsystemslab-uf/DesignSpaceExplorationforSecurity}

 \vspace{-1pt}